\documentclass[11pt]{article}
\usepackage[a4paper,margin=1in]{geometry}
\usepackage{amsmath,amssymb}
\usepackage{booktabs}
\usepackage{hyperref}
\title{GAAE Robustness Evaluation with Cohort-Specific Thresholds\\DELCODE Whole-Brain (Schaefer 400)}
\author{Project Notes}
\date{\today}

\begin{document}
\maketitle

\section*{Objective}
This document describes the robustness protocol used after training a Graph Attention Autoencoder (GAAE) on DELCODE whole-brain functional connectivity data. The method quantifies how stable cohort decisions remain when controlled perturbations are applied to top-confidence test samples.

The pipeline has four core components:
\begin{itemize}
\item reconstruction-error inference,
\item cohort-specific threshold estimation from validation data,
\item top-$k$ test-sample selection per cohort,
\item perturbation-based stability testing.
\end{itemize}

\section*{Data and Model Context}
The method is run after GAAE training/checkpoint loading in:
\begin{itemize}
\item \texttt{MODEL/notebooks/GAAE\_DELCODE\_WHOLE\_BRAIN.ipynb}
\end{itemize}

It evaluates cohorts
\[
\mathcal{C}=\{\text{healthy},\text{ad},\text{mci},\text{converter}\}.
\]

\section*{Per-Sample Reconstruction Error}
For each graph sample $s$, the model outputs reconstructed node features and reconstructed adjacency. The total reconstruction error is defined as
\begin{equation}
E(s)=E_x(s)+\lambda_{\text{adj}}E_{\text{adj}}(s),
\end{equation}
with
\begin{align}
E_x(s) &= \operatorname{MSE}\!\left(\hat{X}_s, X_s\right), \\
E_{\text{adj}}(s) &= \operatorname{BCE}\!\left(\hat{A}_s, A_s\right),
\end{align}
where $\lambda_{\text{adj}}=\texttt{adj\_loss\_weight}$.

Thus, each sample receives a scalar anomaly-like score $E(s)$.

\section*{Cohort-Specific One-vs-Rest Thresholds}
For each cohort $c\in\mathcal{C}$, validation labels are binarized as
\[
y_c=\mathbb{1}[\text{cohort}=c].
\]

Two score orientations are tested:
\begin{align}
\text{high-error score: } &S^{\uparrow}=E, \\
\text{low-error score: } &S^{\downarrow}=-E.
\end{align}

AUC is computed for both orientations:
\[
\mathrm{AUC}^{\uparrow}_c=\mathrm{AUC}(y_c,S^{\uparrow}),
\quad
\mathrm{AUC}^{\downarrow}_c=\mathrm{AUC}(y_c,S^{\downarrow}).
\]

The better orientation is selected:
\[
\text{direction}_c=
\begin{cases}
\text{high}, & \mathrm{AUC}^{\uparrow}_c\ge \mathrm{AUC}^{\downarrow}_c,\\
\text{low}, & \text{otherwise}.
\end{cases}
\]

Using the selected orientation, an ROC curve is built and the threshold is chosen by Youden's index
\begin{equation}
\tau_c=\arg\max_t\left(\mathrm{TPR}_c(t)-\mathrm{FPR}_c(t)\right).
\end{equation}

Mapped back to error space, the cohort threshold $T_c$ is:
\begin{equation}
T_c=
\begin{cases}
\tau_c, & \text{if direction}_c=\text{high},\\
-\tau_c, & \text{if direction}_c=\text{low}.
\end{cases}
\end{equation}

\section*{Cohort Decision Rule}
For any sample with error $E$, cohort-positive prediction for cohort $c$ is
\begin{equation}
\hat{y}_c(E)=
\begin{cases}
\mathbb{1}[E\ge T_c], & \text{if direction}_c=\text{high},\\
\mathbb{1}[E\le T_c], & \text{if direction}_c=\text{low}.
\end{cases}
\end{equation}

This is applied to the test set for every $c\in\mathcal{C}$.

\section*{Top-$k$ Sample Selection per Cohort}
For each cohort $c$, candidates are restricted to true cohort members that are correctly cohort-positive:
\[
\mathcal{S}_c=\{s:\text{cohort}(s)=c\;\wedge\;\hat{y}_c(E(s))=1\}.
\]

A confidence margin is defined by threshold direction:
\begin{equation}
M_c(s)=
\begin{cases}
E(s)-T_c, & \text{if direction}_c=\text{high},\\
T_c-E(s), & \text{if direction}_c=\text{low}.
\end{cases}
\end{equation}

The top $k=5$ samples with largest margin are retained for robustness stress testing.

\section*{Noise and Perturbation Operators}
For each selected sample, perturbations are applied with noise levels
\[
\alpha\in\{0.00,0.05,0.10,0.20,0.30\}
\quad (\text{reported as }0\%,5\%,10\%,20\%,30\%).
\]
Each condition is repeated over $10$ trials.

Methods:
\begin{enumerate}
\item \textbf{feature\_noise}: additive Gaussian perturbation on node features,
\[
X' = X + \epsilon,\quad \epsilon\sim\mathcal{N}(0,\alpha^2\sigma_X^2).
\]
\item \textbf{matrix\_noise\_rebuild}: feature noise as above, then graph is rebuilt by $k$-NN adjacency from perturbed features.
\item \textbf{edge\_perturbation}: random edge deletion plus random edge insertion in proportion $\alpha$.
\item \textbf{conditioning\_noise}: perturb condition vector (age/sex) passed to conditioned decoder.
\end{enumerate}

For each perturbed instance, $E$ is recomputed and transformed into cohort-positive prediction $\hat{y}_c(E)$.

\section*{Stability Metric}
Cohort stability for cohort $c$ under perturbation is defined as
\begin{equation}
\mathrm{Stable}_c = \mathbb{1}[\hat{y}_c(E)=1],
\end{equation}
because selected samples are true members of cohort $c$ and should remain cohort-positive if decisions are stable.

The reported cohort stability rate is
\begin{equation}
\mathrm{CSR}_c(\alpha,m)=\mathbb{E}[\mathrm{Stable}_c\mid \text{noise level }\alpha,\text{ method }m].
\end{equation}

\section*{Plots Produced}
For each cohort $c\in\{\text{healthy},\text{ad},\text{mci},\text{converter}\}$ two plots are generated:
\begin{enumerate}
\item \textbf{Error Drift under Noise}: mean $E$ vs noise level with dashed threshold line at $T_c$.
\item \textbf{Decision Stability under Noise}: $\mathrm{CSR}_c$ vs noise level.
\end{enumerate}

Additional notebook cells can render these two-panel plots independently for each of ad, mci, and converter.

\section*{Saved Artifacts}
Each robustness run stores:
\begin{itemize}
\item \texttt{selected\_top5\_by\_cohort.csv}
\item \texttt{robustness\_summary\_by\_cohort.csv}
\item \texttt{robustness\_details\_by\_cohort.csv}
\item \texttt{cohort\_thresholds.csv}
\item \texttt{robustness\_meta.json}
\end{itemize}

\section*{Interpretation Notes}
\begin{itemize}
\item If mean error drifts across $T_c$ as noise increases, cohort stability should decline.
\item Different cohorts may prefer different threshold directions (high-error or low-error) depending on validation separability.
\item Stability close to $1$ means selected high-confidence cohort members remain cohort-positive under perturbation.
\item Stability collapse at high noise indicates decision boundary sensitivity.
\end{itemize}

\section*{Implementation Reference}
The exact executable implementation is in:
\begin{itemize}
\item \texttt{MODEL/notebooks/GAAE\_DELCODE\_WHOLE\_BRAIN.ipynb} (robustness section and added cohort-specific plotting cells).
\end{itemize}

\end{document}
